% 《大学物理学学习指导》第5版，例16.5—16.7，
% 原书印刷228—229／PDF236—237页；独立证明和边界条件。
\subsection{例题：指数衰变、半衰期与碳-14测年}

\sourceexample{16.5}{228}{236}
\problemtext 设时刻$t=0$有$N_0$个不稳定原子核，
各核种的衰变常数为$\lambda$。
由单位时间内的衰变概率证明
$N(t)=N_0e^{-\lambda t}$。
\workings
独立衰变模型中，短时间$dt$内的衰变概率
为$\lambda dt+o(dt)$。
剩余核数的期望变化满足
\[
 \frac{dN}{dt}=-\lambda N,\qquad
 \frac{dN}{N}=-\lambda\,dt.
\]
积分并代入初始条件$N(0)=N_0$得
\[
 \boxed{N(t)=N_0e^{-\lambda t}}.
\]
令$N(T_{1/2})=N_0/2$，
\[
 \boxed{T_{1/2}=\frac{\ln2}{\lambda}}.
\]
\crosscheck 这里$N(t)$严格说是大量独立随机衰变的
平均存留核数或其宏观近似；有限样本仍有统计起伏。
原书以$\alpha$、$\beta$等衰变作统计示例，
但衰变链补给及时间变化的环境作用会改变单指数模型。

\sourceexample{16.6}{228}{236}
\problemtext 放射性钍-234的半衰期为25天，
求随后一天内核数减少的百分比。
\workings
以天为单位，
$\lambda=(\ln2)/25$，
一天后的比例$N/N_0=e^{-\lambda}=2^{-1/25}$，
因此
\[
 \boxed{f=1-\frac{N}{N_0}
 =1-2^{-1/25}\simeq0.02735=2.73\%}.
\]
\crosscheck 一天衰变率不是恒定的“每日$1/25$”，
应对随时间变化的存留量作指数计算。
原书保留约三位有效数字。

\sourceexample{16.7}{228--229}{236--237}
\problemtext 古生物残留样品的碳-14与碳-12比值
记为$\rho$，空气中对应参考比值为$\rho_0$。
若碳-14半衰期为$T$，证明年代公式。
\workings
设样品死亡后不再与环境交换碳，
且碳-12数目保持不变，则
\[
 \frac{\rho}{\rho_0}=
 \frac{N_{14}(t)}{N_{14}(0)}
 =e^{-\lambda t}.
\]
由$\lambda=\ln2/T$，
取对数得
\[
 \boxed{t=\frac{T}{\ln2}\ln\frac{\rho_0}{\rho}}.
\]
\crosscheck 需要$\rho\le\rho_0$且正值才能对应非负年代。
实际考古测年还需大气碳-14浓度历时变化校准，
以及污染与交换效应的单独检验；
原书题只要求理想模型下的推导。
